\section*{Chapter \ref{ch:g10:price-indices} --- Price Indices and the Cost of Living}
% figdata: figdata/grade-10/12-price-indices/out/ (us_items, us_summary, bren_goods, bren_laspeyres); ledger g10:boskin-1996; Brenland numbers invented.

\begin{solution}{exo:g10:price-indices:1}
$100 \times 2.71/2.00 = 135.5$.
\end{solution}

\begin{solution}{exo:g10:price-indices:2}
Indices 100, $100 \times 9/8 = 112.5$ and $100 \times 10/8 = 125$. Inflation 12.5\,\% in the first year and $125/112.5 - 1 \approx 11.1\,\%$ in the second.
\end{solution}

\begin{solution}{exo:g10:price-indices:3}
College tuition rose the most, to 287.8 (a rise of 188\,\%); computers fell, to 7.6 (a fall of 92.4\,\%). All items rose to 187.0, a rise of 87\,\%.
\end{solution}

\begin{solution}{exo:g10:price-indices:4}
$0.5 \times 120 + 0.3 \times 110 + 0.2 \times 95 = 60 + 33 + 19 = 112$.
\end{solution}

\begin{solution}{exo:g10:price-indices:5}
$2600/1.2 \approx \text{\dal{2166.67}}$ at base-year prices, a real rise of $2166.67/2000 - 1 \approx 8.3\,\%$.
\end{solution}

\begin{solution}{exo:g10:price-indices:6}
$L = 100 \times (3 \times 10 + 5 \times 4)/(2 \times 10 + 5 \times 4) = 100 \times 50/40 = 125$. $P = 100 \times (3 \times 7 + 5 \times 5)/(2 \times 7 + 5 \times 5) = 100 \times 46/39 \approx 117.9$.
\end{solution}

\begin{solution}{exo:g10:price-indices:7}
$0.0553 \times 135.5 + 0.7742 \times 142.9 + 0.0276 \times 130.3 + 0.0737 \times 122.5 + 0.0691 \times 184.0 \approx 7.49 + 110.63 + 3.60 + 9.03 + 12.71 = 143.46$, the same as the
cost ratio up to the rounding of the shares.
\end{solution}

\begin{solution}{exo:g10:price-indices:8}
$100 \times 1.02 \times 1.03 \times 1.025 \approx 107.7$.
\end{solution}

\begin{solution}{exo:g10:price-indices:9}
The \omterm{def:g10:price-indices:laspeyres}{Laspeyres index} rises more: it keeps the old basket, heavy in beef. It overstates the rise in the cost of living, because households can reach a basket they like
nearly as well, with more chicken, for less money than the old basket now costs.
\end{solution}

\begin{solution}{exo:g10:price-indices:10}
(a) $1.06/1.04 - 1 \approx 1.92\,\%$; approximately $6 - 4 = 2\,\%$. (b) $1.20/1.15 - 1 \approx 4.35\,\%$; approximately $20 - 15 = 5\,\%$. The approximation is poor when the
rates are large, since the neglected term grows with their product.
\end{solution}

\begin{solution}{exo:g10:price-indices:11}
$800 \times 132/110 = 960$ dalens.
\end{solution}

\begin{solution}{exo:g10:price-indices:12}
Base year: both goods cost 1 and the household buys 10 of each. Next year good 1 costs 2 and good 2 still 1, and the household buys 12 of good 1 and 8 of good 2. Then $L =
150$ and $P = 100 \times (24 + 8)/(12 + 8) = 160$. Buyers moved towards the good that became dearer. That happens when the price rose because demand rose: a fashion, a
rise in incomes that favours the good, or an improvement in the good that the statistician did not adjust for.
\end{solution}

\begin{solution}{exo:g10:price-indices:13}
The adjusted price is $1100 - 150 = 950$ against 1000: a fall of 5\,\%. Without adjustment, the index would record a rise of 10\,\%.
\end{solution}

\begin{solution}{exo:g10:price-indices:14}
The good enters the basket only at the next update, after its price has already fallen a long way. The fall before entry is never counted, and neither is the gain of the
buyers who could not have it at all before. A basket updated every five years misses up to five years of such falls for each new good.
\end{solution}

\begin{solution}{exo:g10:price-indices:15}
Nominal growth: $393264/248438 \approx 1.583$. With the deflator: $1.583/1.345 \approx 1.177$, real growth of 17.7\,\%, as in \cref{ch:g10:national-accounts}. With the CPI:
$1.583/1.435 \approx 1.103$, only 10.3\,\%. The deflator prices what Brenland produces; the CPI prices what households buy, including imports and excluding investment goods,
so it is the wrong index for the volume of production.
\end{solution}

\begin{solution}{pb:g10:price-indices:1}
Money in dalens.

\textbf{1.} $600 + 8400 + 300 + 800 + 750 = 10850$.

\textbf{2.} $813 + 12001.68 + 391 + 980 + 1380 = 15565.68$.

\textbf{3.} $100 \times 15565.68/10850 \approx 143.46$, which the published CPI rounds to 143.5.

\textbf{4.} Bread 5.53\,\%, rent 77.42\,\%, transport 2.76\,\%, clothing 7.37\,\%, energy 6.91\,\%.

\textbf{5.} Energy rose the most ($0.46/0.25 = 1.84$, $+84\,\%$); clothing the least ($49/40 = 1.225$, $+22.5\,\%$).

\textbf{6.} More bread, transport and clothing; less rent and much less energy. The household bought less of the two goods whose prices rose most (energy and rent) and more of
those whose prices rose least.

\textbf{7.} $840.10 + 11001.54 + 410.55 + 1078 + 920 = 14250.19$.

\textbf{8.} $620 + 7700 + 315 + 880 + 500 = 10015$.

\textbf{9.} $100 \times 14250.19/10015 \approx 142.29$.

\textbf{10.} $143.46 - 142.29 = 1.17$ index points.

\textbf{11.} Laspeyres: $1.4346^{1/10} \approx 1.0368$, 3.68\,\% a year. Paasche: $1.4229^{1/10} \approx 1.0359$, 3.59\,\% a year.

\textbf{12.} About $3.68 - 3.59 \approx 0.09$ percentage points of inflation a year.

\textbf{13.} $30000 \times 1.4346 \approx \text{\dal{43038}}$; on the \omterm{def:g10:price-indices:laspeyres}{Paasche index}, $30000 \times 1.4229 \approx \text{\dal{42687}}$.

\textbf{14.} Yes: the wage grew in the same proportion as the cost of the 2015 basket, so it buys exactly as many such baskets as in 2015. The household is at least as well off,
and better off if it prefers the basket it chose in 2025, which it could afford with money to spare.

\textbf{15.} About 0.09 points a year against the commission's 1.1: the gap measures only substitution between goods, one of its four sources, and in a basket of five goods.

\textbf{16.} $143.5/140.0 = 1.025$: 2.5\,\%.

\textbf{17.} Exactly $1.0462/1.025 - 1 \approx 2.07\,\%$; approximately $4.62 - 2.5 = 2.12\,\%$.

\textbf{18.} The CPI prices households' purchases, including imported goods such as energy, whose prices rose fast; the deflator prices all of Brenland's output, including exports
and investment goods, whose prices rose less.

\textbf{19.} No. The CPI prices a fixed quantity of energy; it does not see that the household now needs less of it. The household's cost of living rose less than the CPI says, which
the \omterm{def:g10:price-indices:laspeyres}{Paasche index} partly captures through the 2025 quantities.

\textbf{20.} Over ten years the Brenland CPI overstated the rise in this household's cost of living by 1.17 index points (143.46 against 142.29), about 0.09 percentage points of
inflation a year.
\end{solution}
